% USERS_MANUAL_EN.tex — xinu-rpi4 User's Manual (English)
% Build: lualatex USERS_MANUAL_EN.tex   (run twice for the ToC)
\documentclass[a4paper,11pt]{article}

\usepackage{fontspec}
\setmonofont{Menlo}[Scale=0.88]   % box-drawing + arrows for code/trees
\usepackage[a4paper,margin=2.2cm]{geometry}
\usepackage{amssymb}
\usepackage{xcolor}
\usepackage{fancyvrb}
\usepackage{longtable}
\usepackage{booktabs}
\usepackage{array}
\usepackage{enumitem}
\usepackage{newunicodechar}
\usepackage[hidelinks]{hyperref}
\usepackage{titlesec}

% --- Unicode glyphs used in body text (code blocks use Menlo directly) ---
\newunicodechar{→}{\textrightarrow}
\newunicodechar{↔}{\ensuremath{\leftrightarrow}}
\newunicodechar{×}{\ensuremath{\times}}
\newunicodechar{…}{\dots}
\newunicodechar{≥}{\ensuremath{\geq}}
\newunicodechar{≠}{\ensuremath{\neq}}
\newunicodechar{✅}{\textcolor{green!55!black}{\checkmark}}
\newunicodechar{⏳}{\textcolor{orange!80!black}{\textbf{\small WIP}}}
\newunicodechar{❌}{\textcolor{red!70!black}{\ensuremath{\times}}}
\newunicodechar{—}{\textemdash}
\newunicodechar{⚠}{\textcolor{orange!80!black}{\textbf{!}}}

% --- code block style ---
\definecolor{codebg}{gray}{0.96}
\DefineVerbatimEnvironment{code}{Verbatim}%
  {fontsize=\small,frame=single,framerule=0.3pt,rulecolor=\color{gray!60},%
   framesep=5pt,xleftmargin=2pt,xrightmargin=2pt}

\setlength{\parindent}{0pt}
\setlength{\parskip}{4pt}
\titlespacing*{\section}{0pt}{14pt}{6pt}
\renewcommand{\arraystretch}{1.25}

\title{\textbf{Xinu \texttt{xinu-rpi4} User's Manual}\\[2pt]
       \large Embedded Xinu for the Raspberry Pi 4 (BCM2711 / Cortex-A72, AArch64)}
\author{}
\date{}

\begin{document}
\maketitle
\vspace{-2.2em}
\begin{center}\rule{0.9\textwidth}{0.4pt}\end{center}

Embedded Xinu port for the Raspberry Pi 4 (BCM2711 / Cortex-A72, AArch64).
The same source tree also cross-builds for the QEMU \texttt{virt} machine.
This manual is written for \textbf{operators} --- for the development
roadmap see \texttt{README.md}.

{\small\tableofcontents}
\bigskip

%======================================================================
\section{Targets and Build Artefacts}

\begin{center}\footnotesize
\begin{tabular}{@{}l l l l@{}}
\toprule
\textbf{Target} & \textbf{Board / machine} & \textbf{Output image} & \textbf{UART0 base}\\
\midrule
\verb|pi4|  & Raspberry Pi 4 (BCM2711) & \verb|kernel8.img|     & \verb|0xFE201000|\\
\verb|qemu| & QEMU \verb|-M virt -cpu cortex-a76| & \verb|kernel_virt.img| & \verb|0x09000000|\\
\bottomrule
\end{tabular}
\end{center}

The load address is \verb|0x80000| for Pi 4 and \verb|0x40080000| for QEMU.
Both images come from one source tree, distinguished by \verb|-D<TARGET>|
macros and a per-target linker script.

\subsection{Feature Matrix (read this first!)}

Not every feature works on every target. Several subsystems run only on
real hardware (Pi 4):

\begin{center}\footnotesize
\begin{tabular}{@{}l c c p{6.6cm}@{}}
\toprule
\textbf{Subsystem} & \textbf{Pi 4} & \textbf{QEMU} & \textbf{Notes}\\
\midrule
PL011 UART0 shell          & ✅ & ✅ & The common substrate\\
Cooperative scheduler (S0) & ✅ & ✅ & \verb|procdemo| / \verb|pingpong|\\
Preemptive scheduler (S1)  & ✅ & ⏳ & 100 Hz timer tick works on Pi 4 (\verb|ticks|)\\
HDMI framebuffer           & ✅ & — & Usable on Pi 4\\
Window manager             & ✅ & — & 1280×960 virtual desktop, Pi 4 only\\
SD card + FAT32            & ✅ & — & EMMC driver is BCM2711-specific\\
Ethernet (GENET)           & ✅ & — & NET-E done (BCM2711 GENET)\\
USB (keyboard / HID)       & ✅ & — & Via VL805 PCIe xHCI (USB-2.0 ports)\\
DHCP / TCP                 & ❌ & — & Code present but dispatch OFF (§7.4)\\
\bottomrule
\end{tabular}
\end{center}

⚠~The source tree contains \texttt{device/genet/}, \texttt{device/sd/},
\texttt{device/usb/xhci/} etc., initialised against \textbf{Pi 4 (BCM2711)
register addresses}. The \verb|-D<BASE>| macros (\verb|GENET_BASE|, \dots) are
defined only for the Pi 4 build; the QEMU build leaves them undefined so the
corresponding code is linked out or no-ops.

\subsection{Boot Sequence (the unusual bits)}

\texttt{loader/boot.S} does more than a textbook bare-metal stub:

\begin{enumerate}[leftmargin=1.6em,itemsep=2pt]
\item \textbf{Mandatory Linux ARM64 Image header.} The Pi 4 EEPROM bootloader
  can refuse to jump into the kernel unless it sees the magic \verb|"ARM\x64"|
  at file offset \verb|0x38|. \texttt{boot.S} places a 64-byte header right at
  \verb|_start:|. Symptom of a missing header: the firmware rainbow pattern
  stays on HDMI and your code never runs.
\item \textbf{Pin secondary cores with MPIDR\_EL1.} The boot core continues;
  cores 1/2/3 enter a \verb|wfe| park loop until S0/S1 wakes them.
\item \textbf{Set up SP using the leex convention.} \verb|sp = _start| (=
  \verb|0x80000|) --- the kernel base address doubles as the initial stack
  pointer.
\item \textbf{Preserve the DTB pointer.} Firmware passes the device tree
  physical address in \verb|x0|; the stub stashes it across BSS clear and
  writes it back to \verb|.data dtb_addr|.
\item \textbf{Zero BSS, then jump to C} (\verb|bl kernel_main|). If
  \verb|kernel_main| ever returns, drop into the same \verb|wfe| park loop.
\end{enumerate}

The stub does \textbf{not} implement an EL2\,\textrightarrow\,EL1 transition
(same assumption holds for QEMU). Adding \verb|armstub=| or \verb|kernel_old=1|
would change this.

%======================================================================
\section{What You Need}

\subsection{Hardware (for a real board)}
\begin{itemize}[leftmargin=1.4em,itemsep=1pt]
\item Raspberry Pi 4 (BCM2711)
\item microSD card with a FAT32 \verb|bootfs| partition (easiest: flash
  Raspberry Pi OS first, then overwrite \verb|kernel8.img| and \verb|config.txt|)
\item 3.3\,V USB-serial adapter, 115200 8N1 --- header pin 8 (TXD\,\textrightarrow\,GPIO14),
  pin 10 (RXD\,\textrightarrow\,GPIO15), and a GND (e.g.\ pin 6)
\item For networking: an Ethernet cable
\end{itemize}

\subsection{Host software}
macOS or Linux, plus an AArch64 cross-toolchain (pick one):
\begin{code}
brew install aarch64-elf-gcc            # GNU
brew install --cask gcc-arm-embedded    # ARM official
\end{code}
\verb|qemu-system-aarch64| for the QEMU target (\verb|brew install qemu|), and a
serial terminal client (\verb|screen|, \verb|minicom|, \verb|picocom|, \dots).

%======================================================================
\section{Build}

Always run \verb|make| from \verb|compile/| --- the Makefile uses \verb|VPATH|
to gather sources from \verb|../loader/|, \verb|../device/...|, etc.

\subsection{Basic commands}
\begin{code}
cd /Users/kodamay/projects/xinu-rpi4/compile

make pi4            # -> compile/kernel8.img       (real Pi 4)
make qemu           # -> compile/kernel_virt.img   (QEMU virt)
make                # = make all = pi4 + qemu
make clean          # remove all objects and images
\end{code}
Objects live in per-variant trees (\verb|obj/pi4/|, \verb|obj/qemu/|), so
building different targets in sequence only recompiles what differs.

\subsection{Toolchain auto-detection}
The Makefile looks for AArch64 GCC as \verb|$(GCCPATH)/bin/aarch64-elf-gcc|
then \verb|aarch64-none-elf-gcc|. Default \verb|GCCPATH| is \verb|/opt/homebrew|:
\begin{code}
make pi4 GCCPATH=$HOME/aarch64/arm-gnu-toolchain-...-aarch64-none-elf
make pi4 CROSS=aarch64-linux-gnu-
\end{code}

\subsection{Key Make variables}
\begin{center}\footnotesize
\begin{tabular}{@{}l l p{7.2cm}@{}}
\toprule
\textbf{Variable} & \textbf{Default} & \textbf{Purpose}\\
\midrule
\verb|GCCPATH| & \verb|/opt/homebrew| & Cross-toolchain install prefix\\
\verb|CROSS|   & (auto-detected)     & Tool prefix (\verb|aarch64-elf-|, \dots)\\
\verb|SDCARD|  & \verb|/Volumes|     & Parent mount point for \verb|make install_*|\\
\verb|DEST|    & \verb|$(SDCARD)/bootfs| & Target FAT32 partition\\
\bottomrule
\end{tabular}
\end{center}

\subsection{Per-target macros (passed via CFLAGS)}
\begin{center}\footnotesize
\begin{tabular}{@{}l l l@{}}
\toprule
\textbf{Macro} & \textbf{Pi 4} & \textbf{QEMU virt}\\
\midrule
\verb|-mcpu|       & \verb|cortex-a72| & \verb|cortex-a76|\\
\verb|UART0_BASE|  & \verb|0xFE201000| & \verb|0x09000000|\\
\verb|MBOX_BASE|   & \verb|0xFE00B880| & (undef)\\
\verb|SD_BASE|     & \verb|0xFE340000| & ---\\
\verb|USB_BASE|    & \verb|0xFE980000| & ---\\
\verb|GIC_BASE|    & \verb|0xFF840000| & ---\\
\verb|PCIE_BASE|   & \verb|0xFD500000| & ---\\
\verb|GENET_BASE|  & \verb|0xFD580000| & ---\\
\verb|HEAP_END|    & \verb|0x40000000| (1G) & \verb|0x50000000|\\
\verb|SKIP_MBOX|   & ---               & defined (simplify)\\
\verb|KERNEL_NAME| & \verb|"kernel8.img"| & \verb|"kernel_virt.img"|\\
\verb|BOARD_NAME|  & \verb|"Pi4"|      & \verb|"virt"|\\
\verb|SOC_NAME|    & \verb|"BCM2711"|  & \verb|"QEMU"|\\
\bottomrule
\end{tabular}
\end{center}

Common CFLAGS:
\verb|-Wall -O2 -ffreestanding -nostdinc -nostdlib -nostartfiles -mgeneral-regs-only|.
The flag keeps most code FPU/NEON-free; only \verb|cc.c| and \verb|llm.c| are
built with FP enabled (CPACR releases the FPU at EL1).

\subsection{Linker scripts}
\begin{center}\footnotesize
\begin{tabular}{@{}l l l@{}}
\toprule
\textbf{Script} & \textbf{Target} & \textbf{Entry / load address}\\
\midrule
\verb|compile/link.ld|      & Pi 4      & \verb|0x80000|\\
\verb|compile/link_virt.ld| & QEMU virt & \verb|0x40080000|\\
\bottomrule
\end{tabular}
\end{center}
Both place \verb|.text.boot| at the start and export \verb|__bss_start| /
\verb|__bss_size| (8-byte units), following the leex convention.

%======================================================================
\section{Writing the SD Card}

\subsection{Pi 4}
\begin{code}
cd compile
make install_pi4 SDCARD=/Volumes
# -> copies kernel8.img and sdcard/config_pi4.txt to /Volumes/bootfs
\end{code}
\verb|make install| is an alias for \verb|install_pi4|.

\subsection{Manual copy (e.g.\ external SD reader)}
\begin{code}
diskutil mount /dev/disk4s1                # -> /Volumes/bootfs
cp compile/kernel8.img /Volumes/bootfs/
cp sdcard/config_pi4.txt /Volumes/bootfs/config.txt
sync
diskutil eject /Volumes/bootfs
\end{code}
⚠~Verify the md5 of \verb|kernel8.img| on the card matches \verb|compile/| ---
booting an old image is the most common time sink.

\subsection{What else must be on the card}
\verb|kernel8.img| + \verb|config.txt| alone won't boot. The same FAT32
partition also needs the Raspberry Pi OS firmware blobs (\verb|bootcode.bin|,
\verb|start4.elf|, \verb|fixup4.dat|, \dots). Fastest path: flash Pi OS first,
then overwrite only \verb|kernel8.img| and \verb|config.txt|.

%======================================================================
\section{Serial Console}
\begin{code}
screen /dev/tty.usbserial-XXXX 115200                       # macOS
minicom -b 115200 -o -D /dev/tty.usbserial-XXXX             # macOS
sudo screen /dev/ttyUSB0 115200                             # Linux
\end{code}
Exit \verb|screen| with \verb|Ctrl-a k| then \verb|y|. Within ~5 seconds of
power-on you should see:
\begin{code}
================================================
  Xinu Pi4 hello (AArch64, BCM2711, kernel8.img)
  PL011 UART0 @ 0xFE201000, 115200 8N1
  bootstrap: leex-style stub + xinu-rpi4 main
================================================

Round 1 phase B/U done -- entering interactive shell.
type `help` for the command list.
xinu-pi4$ _
\end{code}

%======================================================================
\section{Shell Command Reference}
\verb|help| or \verb|?| always shows the live list. The current command set:
\begin{center}\footnotesize
\begin{longtable}{@{}l p{9.0cm}@{}}
\toprule
\textbf{Command} & \textbf{Purpose}\\
\midrule \endhead
\verb|help| / \verb|?|    & List registered commands\\
\verb|echo <words>|       & Echo args back (whitespace-collapsed)\\
\verb|hello|              & Smoke marker --- greeting\\
\verb|mem|                & Show \verb|__bss_start| / \verb|__bss_end| / \verb|_end|\\
\verb|peek <hex_addr>|    & Read a 32-bit MMIO word (e.g.\ \verb|peek 0xfe201018|)\\
\verb|uptime|            & Raw \verb|CNTPCT_EL0| (generic timer counter)\\
\verb|ticks|             & 100 Hz timer tick counter (S1)\\
\verb|ps|                & Core / EL status\\
\verb|halt|              & Mask DAIF + PSCI \verb|SYSTEM_OFF| (QEMU exits cleanly)\\
\verb|reboot|            & Stub --- spins until power-cycle\\
\verb|pwd ls cd mkdir rmdir touch cat write edit rm cp mv tree| & In-RAM filesystem (volatile; §6.3)\\
\verb|cc <file.c>|        & Compile \& run a C program on-device (JIT \textrightarrow{} AArch64; §6.3)\\
\verb|aload| / \verb|amsg| & Load resident AIPL actors / send a message (§6.3)\\
\verb|actordemo| / \verb|selectdemo| & Actor ping-pong / guarded named-message receive (§6.3)\\
\verb|vmtest| / \verb|vmdemand| & VA≠PA remap / demand-paged virtual memory (§6.4)\\
\verb|llm [prompt]|       & Generate text from the baked-in LLM (§6.3)\\
\verb|preempt|           & Demo the timer-driven preemptive round-robin scheduler\\
\verb|pingpong [N]|       & AIPL-style 2-actor cooperative PingPong, N=1..50\\
\verb|procdemo [N]|       & Real 2-process ctxsw demo, N=1..30\\
\verb|usb|               & xHCI / DWC2 USB diagnostics (Pi 4 only)\\
\verb|wifi ...|           & WiFi + mesh: \verb|probe|/\verb|scan|/\verb|up|/\verb|adhoc|/\verb|aodv|/\dots\ (§7.5--§7.6)\\
\verb|rxstat| / \verb|tcpstat| & RX-ring / TCP-listener counters\\
\verb|pan <dx> <dy>| \verb|view| \verb!autopan [on|off]! & Window-manager viewport controls\\
\bottomrule
\end{longtable}
\end{center}

\subsection{procdemo --- real context switch}
\begin{code}
xinu-pi4$ procdemo 3
procdemo: created pid=1 (ping) and pid=2 (pong), iters=3
  [Ping pid=1] tick 1
  [Pong pid=2] tock 1
  [Ping pid=1] tick 2
  [Pong pid=2] tock 2
  [Ping pid=1] tick 3
  [Pong pid=2] tock 3
  [Ping pid=1] exit at iter 3
  [Pong pid=2] exit at iter 3
procdemo: both processes exited; back in shell.
\end{code}
Each \verb|pid=N| is read live from global \verb|currpid|, so the alternation
is direct evidence the scheduler flipped contexts. Stacks are not reclaimed on
\verb|proc_exit| yet --- S1 (clock IRQ) will let the dispatcher reap them.

\subsection{pingpong vs procdemo}
\begin{center}\footnotesize
\begin{tabular}{@{}l p{4.4cm} p{5.0cm}@{}}
\toprule
\textbf{Aspect} & \verb|pingpong| & \verb|procdemo|\\
\midrule
Actors     & Two static (\verb|Ping|, \verb|Pong|) & Real procs via \verb|proc_create|\\
Switch     & Single-stack dispatcher loop & Real ctxsw via \verb|ctxsw.S|\\
Terminate  & Both inboxes empty & Both call \verb|proc_exit()|\\
\bottomrule
\end{tabular}
\end{center}

\subsection{On-device tooling (filesystem, C JIT, actors, LLM)}
Beyond the demos above, the shell carries a self-contained toolchain --- you can
write, compile, and run code on the board itself.
\begin{itemize}[leftmargin=1.4em,itemsep=1pt]
\item \textbf{In-RAM filesystem} --- \verb|pwd| \verb|ls| \verb|cd| \verb|mkdir|
  \verb|rmdir| \verb|touch| \verb|cat| \verb|write| \verb|edit| \verb|rm|
  \verb|cp| \verb|mv| \verb|tree| operate on a small in-memory tree (volatile;
  cleared on reboot).
\item \textbf{C JIT (\texttt{cc})} --- \verb|cc <file.c>| compiles a C subset to
  native AArch64 and runs it in place; the same compiler is reachable over HTTP
  as \verb|POST /compile| (body = C source).
\item \textbf{AIPL actors} --- \verb|aload <file.c>| loads resident actors;
  \verb|amsg <actor> <method> [arg]| sends a message. \verb|actordemo| runs a
  2-actor ping-pong as real Xinu processes; \verb|selectdemo| shows a guarded
  receive. HTTP \verb|/actor|, \verb|/send|, \verb|/gc| expose the actor
  inventory, message send, and the actor-pool GC.
\item \textbf{Embedded LLM (\texttt{llm})} --- a tiny transformer is baked into
  the image; \verb|llm [prompt]| generates text on-device (also \verb|/chat|).
\end{itemize}

\subsection{Virtual memory (\texttt{vmtest} / \texttt{vmdemand})}
The kernel runs with the MMU on (identity map) plus a \textbf{demand-paged
window} at VA \verb|0x80000000|..\verb|0x80400000| (4 MiB, 512-frame pool):
\begin{itemize}[leftmargin=1.4em,itemsep=1pt]
\item \verb|vmtest| --- map a physical page at a different virtual address and
  prove the translation (VA ≠ PA).
\item \verb|vmdemand| --- touch 64 pages in the demand window; the first run
  takes \verb|0x40| page faults (one per page) and reads back OK, the second run
  takes \verb|0x0| faults (already mapped). Check the counter with \verb|/fault|.
\end{itemize}
Both run locally or remotely, e.g.\
\verb|curl "http://192.168.3.100/shell?cmd=vmdemand"|.

%======================================================================
\section{Networking (Pi 4)}
The network stack runs on the Pi 4 GENET (BCM2711).

\textbf{What works:} ARP request/reply; ICMP echo (ping); static IP/MAC
\verb|192.168.3.100| / \verb|d8:3a:dd:a7:fd:bf| (set in \verb|loader/main.c:707|
area); raw broadcast/unicast TX; 16-slot RX ring.

\subsection{Ping from a Mac}
\begin{code}
sudo arp -s 192.168.3.100 d8:3a:dd:a7:fd:bf
ping 192.168.3.100
\end{code}
RTT is ~900\,ms (rate-limited by the window-manager tick). Quick shell checks:
\verb|rxstat| and \verb|ticks|.

\subsection{DHCP / TCP status}
\verb|system/dhcp_client.c| and \verb|system/tcp_server.c| ship in the tree but
their \verb|rx_tick| dispatch is \textbf{OFF}: repeated DHCP broadcast TX
degrades the GENET TX ring and stalls ICMP, and routing \verb|tcp_handle_packet()|
from \verb|rx_tick| stops subsequent ICMP echo replies. Re-enabling is planned in
a future sprint (NET-G bisect).

\subsection{WiFi (BCM43455)}
The Pi 4 has an on-board BCM43455 WiFi chip (separate from the wired GENET
above). It is driven entirely from the serial shell (\verb|xinu-pi4$|): bring up
the firmware, scan, then connect. It does \textbf{not} auto-connect at boot.
\begin{center}\footnotesize
\begin{tabular}{@{}l p{8.6cm}@{}}
\toprule \textbf{Command} & \textbf{What it does}\\ \midrule
\verb|wifi probe| & M0/M1 bring-up: download the firmware to the chip (once per boot)\\
\verb|wifi scan| & escan for nearby APs\\
\verb|wifi up <ssid> <pass>| & join (WPA2-PSK) + DHCP; starts the persistent ARP/ICMP responder\\
\verb|wifi join <ssid> <pass>| & join only (no DHCP)\\
\verb|wifi dhcp| & request a DHCP lease\\
\verb|wifi off| & radio down / disconnect\\
\verb|wifi ping <a.b.c.d> [count]| & ICMP echo client\\
\verb|wifi serve [secs]| & run the ARP/ICMP responder for N seconds (default 30)\\
\verb|wifi resolve <host>| / \verb|wifi web <host>| & DNS / DNS + HTTP GET\\
\verb|wifi ntp [a.b.c.d]| & NTP time client\\
\verb|wifi tftp <ip> <file>| / \verb|wifi netboot <ip> <file>| & TFTP fetch / fetch + chainload\\
\bottomrule
\end{tabular}
\end{center}
\begin{code}
xinu-pi4$ wifi probe          # download firmware (once per boot)
xinu-pi4$ wifi scan           # see what's around
xinu-pi4$ wifi up MyHome-5G mypassword
wifi up: connected; ARP/ICMP responder is now persistent.
xinu-pi4$ wifi ping 8.8.8.8
\end{code}
WiFi does not auto-reconnect after a reboot --- re-run \verb|wifi probe| +
\verb|wifi up| each boot. \verb|wifi netboot <ip> <file>| fetches a kernel over
WiFi and chainloads it (a no-SD-swap update path, like the wired chainload).

\subsection{Mesh networking: multiple Xinu nodes (MANET ad-hoc)}
Several Pi 4 (and Pi 3) Xinu boards can form a peer-to-peer mesh with \textbf{no
access point}, using 802.11 IBSS (ad-hoc) mode plus on-demand AODV routing ---
the same MANET stack used in the drone-HIL demo (Pi 4 + Pi 3 nodes). Each node
joins the same ad-hoc cell with a distinct node number:
\begin{code}
xinu-pi4$ wifi probe                       # once per boot
xinu-pi4$ wifi adhoc <cell-ssid> [ch] [node]
\end{code}
\begin{itemize}[leftmargin=1.4em,itemsep=1pt]
\item \verb|<cell-ssid>| --- the cell name; \textbf{all nodes must use the same
  name and channel}.
\item \verb|[ch]| --- 802.11 channel (default 6).
\item \verb|[node]| --- this node's number (default 1); it gets the static IP
  \verb|10.0.0.<node>/24|. \textbf{Give every node a distinct number.}
\end{itemize}
Example --- a 3-node cell on channel 6:
\begin{code}
xinu-pi4$ wifi adhoc mesh1 6 1      # -> 10.0.0.1
xinu-pi4$ wifi adhoc mesh1 6 2      # -> 10.0.0.2
xinu-pi4$ wifi adhoc mesh1 6 3      # -> 10.0.0.3
\end{code}
Nodes within radio range reach each other directly at \verb|10.0.0.x|
(\verb|wifi ping 10.0.0.2|).

\textbf{Multi-hop (AODV).} When a destination is not in direct range, discover a
route on demand with \verb|wifi aodv <ip>|:
\begin{code}
xinu-pi4$ wifi aodv 10.0.0.3
[wifi] === AODV discover 10.0.0.3 ===
[wifi] aodv: RREQ id=1 for 10.0.0.3
[wifi] *** AODV route: 10.0.0.3 via 10.0.0.2, 2 hop ***
\end{code}
The minimal AODV module (M13, RFC 3561 core) broadcasts an RREQ (UDP port 654);
each in-range node relays it and records a reverse route; the destination answers
with an RREP that installs the forward route. Every node relays for its peers
automatically (route table: up to 16 entries). IBSS / ad-hoc is independent of
the infrastructure \verb|wifi up| mode (run \verb|wifi off| to leave an AP first)
and is not restored after reboot --- re-run \verb|wifi probe| + \verb|wifi adhoc|
on each node.

\subsection{HTTP control plane \& remote shell}
\verb|system/tcp_server.c| runs an HTTP server on the wired interface (default
port 80, \verb|192.168.3.100|), so the board can be driven entirely over the
network:
\begin{center}\footnotesize
\begin{tabular}{@{}l p{8.4cm}@{}}
\toprule \textbf{Route} & \textbf{What it does}\\ \midrule
\verb|GET /shell?cmd=<cmd>| & run any shell command, return its output (no stdin)\\
\verb|POST /compile| & JIT-compile and run a C program (body = source)\\
\verb|GET /chat| & on-device LLM chat\\
\verb|GET /actor| \verb|/send| \verb|/gc| \verb|/jitstats| & actor inventory / send / GC / JIT counters\\
\verb|GET /usbdiag| \verb|/pcie-init| \verb|/pcie-enum| \verb|/xhci-reset| & USB / PCIe bring-up diagnostics\\
\verb|GET /fault| \verb|/mmio-read| \verb|/mmio-write| \verb|/mmio-sweep| & fault counters + raw MMIO peek/poke\\
\verb|POST /reboot| & BCM2711 watchdog reset\\
\verb|POST /chainload| & upload + jump to a new kernel (no SD swap)\\
\bottomrule
\end{tabular}
\end{center}
Runtime diagnostics use these HTTP counters because \verb|uart_puts| deadlocks in
the HTTP-worker context --- all PCIe/xHCI bring-up logging is done at boot over
serial.

\textbf{Network kernel update (chainload --- no SD swap).} \verb|POST /chainload|
uploads a new kernel and jumps to it in RAM. A helper script wraps the upload:
\begin{code}
python3 tools/remote_chainload.py 192.168.3.100 compile/kernel8.img
\end{code}
A bad image just needs a power cycle (the SD card is untouched), which keeps the
dev loop fast. It needs the HTTP server to be responsive to accept the upload.

%======================================================================
\section{Running under QEMU}
\begin{code}
cd compile
make qemu          # interactive -- Ctrl-A X to quit
make qemu-smoke    # canned input, writes qemu-smoke.log
\end{code}
Under QEMU: \verb|MIDR_EL1 = 0x414fd0b1| (QEMU's published Cortex-A76 part
number; real Pi 4 hardware runs Cortex-A72); \verb|halt| cleanly exits because
\verb|virt| handles PSCI \verb|SYSTEM_OFF|; no networking, USB, SD, or HDMI.

%======================================================================
\section{Important config.txt Lines}
\verb|sdcard/config_pi4.txt| (copied to the card's \verb|config.txt| by
\verb|make install_pi4|):
\begin{center}\footnotesize
\begin{tabular}{@{}l p{7.6cm}@{}}
\toprule
\textbf{Line} & \textbf{Meaning}\\
\midrule
\verb|arm_64bit=1|              & Boot in AArch64\\
\verb|kernel=kernel8.img|       & Pi 4 kernel name\\
\verb|kernel_address=0x80000|   & Load address\\
\verb|enable_uart=1|            & Enable UART\\
\verb|uart_2ndstage=1|          & Print second-stage bootloader log on UART\\
\verb|dtparam=uart0=on|         & Route UART0 to GPIO14/15\\
\verb|init_uart_clock=48000000| & Pin the PL011 reference clock to 48 MHz\\
\verb|hdmi_force_hotplug=1|     & Force HDMI hotplug\\
\verb|framebuffer_width=640|    & Width of firmware-allocated framebuffer\\
\verb|framebuffer_height=480|   & Height\\
\verb|framebuffer_depth=32|     & bpp\\
\bottomrule
\end{tabular}
\end{center}
The KMS overlay (\verb|dtoverlay=vc4-kms-v3d|) is left in --- Xinu still
consumes the firmware's simple-framebuffer via the VC mailbox.

%======================================================================
\section{USB / Input Device Support}
The Pi 4's USB-A ports go through \textbf{VL805 (PCIe xHCI)}; a hand-rolled
xHCI driver brings up USB keyboards and mice.
\begin{itemize}[leftmargin=1.4em,itemsep=1pt]
\item \textbf{Use the USB-2.0 (black) ports.} They route through the USB-2.0 hub
  (Genesys, working TT), so low-speed interrupt-IN split transactions are
  delivered. The USB-3.0 (blue) ports sit behind the VIA USB3 hub whose path
  does not deliver the periodic transfer (the cursor won't move).
\item The shell's \verb|usb| command is a diagnostic dump for the DWC2 (USB-C
  OTG) HCD, not a general device-control interface.
\item VL805 firmware leaves the PCIe RC in reset on an SD boot, so boot issues a
  \verb|NOTIFY_XHCI_RESET| mailbox tag (\verb|0x00030058|, devid \verb|0x100000|)
  to bring xHCI up (details in \verb|NEXT_SESSION_PI4.md|).
\end{itemize}

\begin{center}\footnotesize
\begin{tabular}{@{}l l l@{}}
\toprule
\textbf{Board} & \textbf{Input} & \textbf{Output}\\
\midrule
Pi 4  & USB-A keyboard / UART0 serial & UART0 serial + HDMI framebuffer\\
QEMU  & qemu stdio (serial)          & qemu stdio (serial)\\
\bottomrule
\end{tabular}
\end{center}

%======================================================================
\section{Troubleshooting}
\begin{center}\footnotesize
\begin{tabular}{@{}p{5.2cm} p{8.2cm}@{}}
\toprule
\textbf{Symptom} & \textbf{What to check}\\
\midrule
No output on serial            & device name + speed in \verb|screen|; TX/RX swapped; GND\\
Banner appears, then stops     & Old \verb|kernel8.img| on SD --- rebuild and recopy\\
No echo at the prompt          & Local echo on the host terminal --- keep it OFF\\
\verb|peek| returns 0xFFFFFFFF  & Address not mapped; MMU is flat ID\\
USB keyboard/mouse dead        & Make sure it's on a USB-2.0 (black) port (§11)\\
Ping fails on Pi 4             & \verb|sudo arp -s ...| done?; cable; \verb|rxstat| increments?\\
\verb|make qemu| fails          & Check \verb|qemu-system-aarch64| version ($\geq$11)\\
\verb|make install| permission  & It uses \verb|sudo cp| --- provide the sudo password\\
Reboot loop on power-on        & Does \verb|kernel=| match the file on the card?\\
Rainbow pattern stays on HDMI  & Likely missing ARM64 Image header --- rebuild\\
\bottomrule
\end{tabular}
\end{center}
Generic recovery flow:
\begin{code}
cd compile && make clean && make pi4
make install_pi4 SDCARD=/Volumes
diskutil eject /Volumes/bootfs
# Reseat the SD card and power-cycle the board.
\end{code}

%======================================================================
\section{Directory Layout (operator view)}
\begin{code}
xinu-rpi4/
├── compile/                # run `make` here
│   ├── Makefile
│   ├── kernel8.img         # Pi 4 (built by `make pi4`)
│   └── kernel_virt.img     # QEMU (built by `make qemu`)
├── sdcard/
│   └── config_pi4.txt      # Pi 4 (copied by install_pi4)
├── README.md               # Developer-facing roadmap + internals
├── USERS_MANUAL_EN.md      # Markdown source of this manual
├── USERS_MANUAL_JA.md      # Japanese version
└── NEXT_SESSION_PI4.md     # Handoff notes for ongoing work
\end{code}

%======================================================================
\section{Recipe Book}
\textbf{From a fresh checkout to a running Pi 4:}
\begin{code}
cd /Users/kodamay/projects/xinu-rpi4/compile
make clean && make pi4
make install_pi4 SDCARD=/Volumes
diskutil eject /Volumes/bootfs
screen /dev/tty.usbserial-XXXX 115200
\end{code}
\textbf{Quick command sanity check (QEMU):}
\begin{code}
cd compile && make qemu
xinu-pi4$ procdemo 5
xinu-pi4$ pingpong 3
xinu-pi4$ halt          # Ctrl-A X to exit QEMU
\end{code}

%======================================================================
\section{References}
\begin{itemize}[leftmargin=1.4em,itemsep=1pt]
\item Upstream Xinu (32-bit): \url{https://github.com/yaskodama/xinu-rpi}
\item AArch64 boot stub origin: \url{https://github.com/radlyeel/leex}
\item Shell centry pattern origin: \url{https://github.com/davidxyz/xinuPi}
\item Development roadmap (Round 1 plan): under
  \verb|aice-pi-evolution/experiments/| in the abclcp-project repo
\end{itemize}

\section{License}
Inherits the BSD-style licenses of upstream Xinu and leex. See \verb|LICENSE|
once the source-of-truth file is added.

\end{document}
